\section{Positive collision series and spectral bins}\label{tra:collision-series}

A regular-rank bound does not retain the distribution of moment mass across diagram levels. The positive series below does: its coefficients include the tail-tableau multiplicities, and Cauchy bounds can extract individual level contributions.

A tuple collision calculation retains a positive sum of traces at each representation level. The weight of a representation in this sum is the dimension of the diagram below its first row. Branching extracts these weighted sums from the exponential generating function for tuple traces. We estimate the resulting coefficients through the collision graph of a single realized sweep. Positivity then permits a high-power estimate before the regular multiplicities are restored.

For the dense trace, the input is a compatibility bound under independent component laws with specified supremum densities. We divide positive subgroup operators into spectral bins. Each bin has a controlled rank and therefore a controlled supremum density after its squared coefficient kernel is normalized. This gives the weak trace estimate to which the positive collision coefficients will be applied.

Let $n=2^d$, $d\ge1$. Write $T_n$ for a sweep and $\mathcal Q_n=T_nT_n^*$, acting on the regular representation. One sweep corresponds to $d$ physical shuffles by Section~\ref{sec:prelim}. Products act right-to-left, as in Section~\ref{sec:prelim}; a row move followed by a column move therefore places the column operator on the left. Traces are unnormalized.

\subsection{Supremum-density compatibility bounds}
\begin{proposition}\label{tra:collision-series:compatibility}
\label{tra:collision-series:compatibility-statement}
The following estimate holds for sufficiently large \(n\) factored as \(n=l m\) with \(1/2\le l/m\le 2\). Let \(G\) be the permutation group of \([l]\times[m]\), with \(H\) its subgroup of permutations acting separately in each row and \(K\) the subgroup acting separately in each column. For \(h\in H,\ k\in K\), consider the event
\begin{equation}
hk\in KH. \label{tra:collision-series:eq2}
\end{equation}
Suppose \(h,k\) are chosen with independent component permutations (all \(l+m\) permutations independent). Use \(L_v\) as bounds on the supremum densities, relative to uniform, for the individual permutation laws, indexing these components by \(v\). Then for an absolute \(C\),
\begin{equation}
\mathbb P\{hk\in KH\}
\ \le\
\exp\left(-n+C n^{0.55}+\frac{C}{\log n}\sum_v\log L_v\right). \label{tra:collision-series:eq3}
\end{equation}
\end{proposition}
\begin{proof}\label{tra:collision-series:compatibility-proof}
\label{tra:collision-series:bundle-packing-and-entropy}
To prove this, first invert the permutations: \(hk\in KH\) if and only if \(k^{-1}h^{-1}\in HK\). The latter event concerns a row move \(h^{-1}\) followed by a column move \(k^{-1}\). Inversion preserves independence of the component laws and their supremum density bounds. We describe this event as a constraint on two arrays. For intermediate cell \((i,t)\) (after \(h^{-1}\) and before \(k^{-1}\)), let \(J_i(t)\) be its column of origin and \(R_t(i)\) its row of destination. Thus \(J_i\), \(i\in[l]\), are permutations of \([m]\), and \(R_t\), \(t\in[m]\), are permutations of \([l]\). Constraint \eqref{tra:collision-series:eq2} says that the pairs
\[
\big(J_i(t),R_t(i)\big),\qquad (i,t)\in[l]\times[m],
\]
are all distinct. Indeed this is necessary to be obtainable with a column layer first and a row layer second; if it holds, the column layer can first give every point its intended destination row, after which a row layer can give the intended columns. Here by row or column layer we mean a permutation in \(H\) or \(K\), respectively. Call arrays satisfying the constraint valid.

The random-order exposure below is related to Radhakrishnan's entropy proof of Br\'egman's theorem~\cite{Radhakrishnan1997} and its higher-dimensional development by Linial and Luria~\cite[Sections~2.2--3.2]{LinialLuria2014}. Here the conditional predictions need not be uniform: the estimate retains the component entropy deficits and separately charges concentrated entries.

Use now an arbitrary distribution \(\nu\) supported on valid arrays; entropies, denoted by \(\mathsf h\), are in nats. Write \(D_\mathrm{u}\) for relative entropy to uniform on permutations (or product uniform when applied to a group of array components), using the marginal under \(\nu\) in this notation. We will show
\begin{equation}
\begin{split}
\mathsf h(J,R) \ \le\ 
 &\sum_i\mathsf h(J_i)+\sum_t\mathsf h(R_t)
 - n+C n^{0.55}\\
 &\hspace{4mm}+\frac{C}{\log n}\left(\sum_i D_\mathrm{u}(J_i)+\sum_t D_\mathrm{u}(R_t)\right). 
\end{split}\label{tra:collision-series:eq4}
\end{equation}

Fix \(\eta=1/100\), \(B=\lceil n^\eta\rceil\), and write
\[
b_{it}=\#\{t': R_{t'}(i)=R_t(i)\},\qquad W=\#\{(i,t): b_{it}\ge B\}.
\]
There is an absolute \(c>0\) (we can use \(c=\eta/4\)) such that
\begin{equation}
\mathbb E_\nu W\ \le\ \frac{D_\mathrm{u}(R)+\log 2}{c\log n}. \label{tra:collision-series:eq5}
\end{equation}
In fact, consider first independent uniform column permutations. A bundle consists of a row \(i\), a row value \(r\), and \(B\) columns, with a prescription that all those columns have \(R_t(i)=r\). A disjoint family of bundles means the prescribed cells do not overlap. For any such family of size \(s\), the probability of satisfying it is at most \((e/l)^{Bs}\). Indeed, each column with \(a\) cells specified gives either probability zero or \(1/(l)_a\), where \((l)_a\) denotes a falling factorial, and \((l)_a\ge (l/e)^a\). For any outcome we can pack bundles it satisfies so as to have \(2B\) times their number at least \(W\): use disjoint bundles for each row and value with count at least \(B\). Thus \(n^{c W}\) is at most the sum of \(n^{2cBs}\) over satisfied disjoint families, including the empty family. Its expectation under product uniform is bounded by
\[
\sum_{s\ge 0}\left(l^2 \binom mB (e/l)^B n^{2cB}\right)^s\ \le\ 2
\]
for large \(n\), by \(m/l\le 2\). This moment bound gives \eqref{tra:collision-series:eq5} by the relative entropy inequality (or changing measure and using Jensen's inequality).

The bundle estimate charges repeated destination values to the joint entropy deficit of the column array. We now expose the row array to obtain one nat of compatibility saving per cell, except for these repetitions and for heavy conditional atoms.

We bound the conditional entropy for \(J\) given \(R\) by revealing its rows sequenced by independent uniform priorities in \([0,1]\), independent of the arrays. Use decreasing priority, so given the current row's priority \(x\), each other row has probability \(x\) of being later. Within a row reveal cells in fixed order. For cell \((i,t)\), use \(p_j\) to denote the conditional probabilities for its value \(J_i(t)=j\) under the marginal law of \(J_i\), given the previously revealed part of just that row. There are \(u\) remaining columns to fill in the row, so \(p\) is supported on \(u\) remaining values (allowed to have zero probabilities there). The sum of expected \(-\log p_j\) for the true values over cells gives \(\sum_i\mathsf h(J_i)\).

We give the savings in the entropy upper bound arising from earlier rows. Call a candidate \(j\) heavy if \(p_j>B/u\), and otherwise light. Let their class masses be \(P_\mathrm{heavy},P_\mathrm{light}\). For light candidates, consider availability, meaning that the pair with this \(j\) and current \(R_t(i)\) has not appeared in earlier rows. Use as test probabilities for the entropy bound \(p_j\) unchanged for heavy candidates; for light candidates redistribute their class mass on available light candidates proportionally to \(p_j\). If no such positive mass is available, omit the light class. These probabilities (possibly a subprobability) use only the permitted history with \(R\) known, and the true value has positive test probability almost surely. By the entropy or cross-entropy inequality at every step, this gives the upper bound on conditional entropy equal to the base terms from \(-\log p_j\), plus terms of the following form:
\begin{itemize}
\item a contribution zero if the true value is heavy;
\item for a light true value, the logarithm of the fraction of \(P_\mathrm{light}\) available in positive-probability light candidates.

\end{itemize}
Here and below it is the expectation of such contributions that enters the bound.

The cell gives an expected contribution of about \(-1\) outside of some losses, as follows. Fix the full arrays first, and suppose the true value is light, \(P_\mathrm{light}\ge 1/2\), and \(b_{it}<B\). Every pair in question occurs exactly once in the valid arrays. The total fraction in the light class from pairs located within row \(i\) is at most \(2B^2/u\): there are at most \(b_{it}\) of them and each light mass is at most \(B/u\). For pairs outside this row the probability of availability is \(x\), given the current priority. Hence the average log fraction, bounding by Jensen's inequality and integrating priority, is at most
\[
\int_0^1\log(x+2B^2/u)\,dx
\ \le\ -1+C'\frac{B^2}{u}\log n
\]
with an absolute \(C'\); this follows by integration, for large \(n\), also when \(2B^2/u>1\). We have used the independence of the row priorities from the full arrays.

Always, the contribution is at most zero. The losses (missed savings of 1) on summing from \(b_{it}\ge B\) can thus be charged to \(\mathbb E_\nu W\); the total losses from heavy true value or \(P_\mathrm{light}<1/2\) are bounded by 3 times the sum of expected \(P_\mathrm{heavy}\). Conditionally using just the row history, heavy mass is bounded by
\[
P_\mathrm{heavy}\ \le\ \frac{D(p\|\operatorname{Unif}_u)}{\log B-1}
\]
for large \(n\), using \(\operatorname{Unif}_u\) on the remaining values. To see this use \(z\log z-z+1\ge 0\) with \(z=u p_j\) in the divergence, and bound it below by \(z(\log B-1)\) for heavy values. The sum of these expected conditional divergences (the numerators) is \(\sum_i D_\mathrm{u}(J_i)\).

Applying the entropy chain rule (with independent priorities) now gives
\[
\mathsf h(J\mid R)\ \le\ \sum_i\mathsf h(J_i)-n
  +C'\sum_i\sum_{u=1}^m (B^2/u)\log n
  +\mathbb E_\nu W
  +\frac{3}{\log B-1}\sum_i D_\mathrm{u}(J_i).
\]
The double sum error with its prefactor is at most \(C'' n^{0.55}\), as \(l=O(\sqrt n)\) and \(B=\lceil n^{1/100}\rceil\). In adding \(\mathsf h(R)\), use
\[
\mathcal C_R=\sum_t\mathsf h(R_t)-\mathsf h(R)\ \ge\ 0,\qquad
D_\mathrm{u}(R)=\sum_t D_\mathrm{u}(R_t)+\mathcal C_R.
\]
Then \eqref{tra:collision-series:eq5} shows we get \eqref{tra:collision-series:eq4} (the positive error coming from \(\mathcal C_R/(c\log n)\) is absorbed by the saving of \(\mathcal C_R\)). This proves the entropy claim.

To conclude \eqref{tra:collision-series:eq3}, take \(\nu\) to be the law of the independent components conditioned on validity, assuming positive probability. Write \(\rho=\bigotimes_v \rho_v\) for the independent law before conditioning, and \(\nu_v\) for the marginal at \(v\). In passing between permutation components and the arrays we may have taken inverses, which does not affect the given supremum bounds. Minus the log probability of validity equals
\[
D(\nu\|\rho)
 = \sum_v D(\nu_v\|\rho_v)+ \sum_v \mathsf h(\nu_v)-\mathsf h(\nu).
\]
In \eqref{tra:collision-series:eq4}, the individual deficits \(D_\mathrm{u}\) are at most \(D(\nu_v\|\rho_v)+\log L_v\). Using \eqref{tra:collision-series:eq4} and \(C/\log n\le 1\) for sufficiently large \(n\), we get \eqref{tra:collision-series:eq3}.
\end{proof}
\subsection{Spectral bins in the fourth trace}
\label{tra:collision-series:balanced-split}
The compatibility estimate concerns probability laws, whereas subgroup coefficient kernels may be signed. Squaring a coefficient kernel produces a nonnegative law after normalization, and splitting first into spectral bins gives the density control the estimate requires. From now on, split \(d\) into \(\lfloor d/2\rfloor+\lceil d/2\rceil\). Use
\[
m=2^{\lfloor d/2\rfloor},\qquad l=2^{\lceil d/2\rceil},
\]
and take \(d\) large enough that both are at least 2. The first part of the pass works independently within the \(l\) rows of size \(m\) and the second works independently within the \(m\) columns of size \(l\). They are products of the corresponding smaller passes on these subgroups, extended in convolution action to \(G\). Denote them by \(F_H,F_K\), so \(T_n=F_K F_H\).
\begin{proposition}\label{tra:collision-series:weak-moment}
\label{tra:collision-series:weak-moment-statement}
Suppose the two smaller sizes satisfy
\[
 \operatorname{Tr}\mathcal Q_m^{q_m}\le 1+\tfrac18,
 \qquad
 \operatorname{Tr}\mathcal Q_l^{q_l}\le 1+\tfrac18,
 \qquad q_m,q_l\ge300.
\]
No uniform upper bound on these powers is needed at this step. These are the child bounds for the invariant \eqref{tra:collision-series:eq1} closed below. There is an absolute \(A\) such that with
\[
p=(1+A/\log n)\max(q_m,q_l)
\]
we have
\begin{equation}
\operatorname{Tr} \mathcal Q_n^p\ \le\ \exp(n^{0.57}) \label{tra:collision-series:eq6}
\end{equation}
for all sufficiently large \(n\).
\end{proposition}
\begin{proof}\label{tra:collision-series:weak-moment-proof}
\label{tra:collision-series:spectral-bin-recursion}
Set \(X=F_H F_H^*,\ Y=F_K^*F_K\), positive operators. Since \(T_n^*T_n=F_H^*YF_H\), the positive operators \(T_n^*T_n\), \(\mathcal Q_n\), and \(Y^{1/2}XY^{1/2}\) have the same eigenvalues, including zeros. Their \(p\)-th powers therefore have the same trace. Lemma~\ref{lem:positive-power-comparison}, with exponent $p\ge2$,
gives
\[
\operatorname{Tr} \mathcal Q_n^p \ \le\ \operatorname{Tr}\left(X^{p/2} Y^{p/2} X^{p/2} Y^{p/2}\right).
\]
Use \(x(h),y(k)\) for the kernels of \(X^{p/2}, Y^{p/2}\), respectively, normalized on their own groups, meaning that \(X^{p/2}\) acts by convolution with \(x(h)/|H|\), and similarly for \(Y^{p/2}\). The powers here indeed come from these groups, as extension to \(G\) acts by the same matrix blocks on corresponding cosets. Both \(x\) and \(y\) factor as products over permutation components in their respective groups. Apply Lemma~\ref{lem:coefficient-compatibility} to the positive
line powers composing $X^{p/2}$ and $Y^{p/2}$. In the present
orientation compatibility is $hk\in KH$; the lemma gives the
unnormalized squared-coefficient form
\begin{equation}
\operatorname{Tr} \mathcal Q_n^p \ \le\
\frac{|G|}{|H||K|}
\mathbb E_{H,K}\!\left[{\boldsymbol 1}_{hk\in KH}\, |x(h)y(k)|^2\right], \label{tra:collision-series:eq7}
\end{equation}
where expectation is independent uniform. This is where we can apply the grid estimate, but the kernels should first be split into sizes.

Each permutation component \(v\) of \(H,K\) uses a group of degree \(a=m\) or \(l\). Its kernel factor in \(x\) or \(y\) is for \(D_v^{p/2}\), where \(D_v\) is a copy of \(T_a^*T_a\) or \(T_a T_a^*\) on the regular representation of the component. Split according to spectral bins indexed by \(b=0,1,\ldots\), grouping eigenvalues of \(D_v^{q_a}\) in
\[
\big(e^{-(b+1)}, e^{-b}\big].
\]
Only nonzero eigenvalues matter. These pieces, also of convolution form by spectral calculus, have ranks \(r_{v,b}\le 2 e^{b+1}\) by the input trace bound. Write \(z_{v,b}\) for the normalized kernel for \(D_v^{p/2}\) restricted to the bin, zero on the orthogonal complement. For expectation uniform on the component group, using the stated child trace bounds, we have
\[
w_{v,b}:=\mathbb E |z_{v,b}|^2\ \le\
r_{v,b}\exp(-b p/q_a)\ \le\ 2 e\,\exp\left(-b A/\log n\right).
\]
In the first inequality we used that uniform mean square equals the squared Hilbert-Schmidt norm (ordinary, for the convolution matrix in the regular representation). Also, for each nonempty bin,
\[
\frac{\sup |z_{v,b}|^2}{w_{v,b}}\ \le\ r_{v,b}.
\]
This is the rank bound for a positive coefficient density in
Lemma~\ref{lem:coefficient-compatibility}, applied to the bin operator.

Each nonempty bin now supplies a probability density $|z_{v,b}|^2/w_{v,b}$ bounded by its rank. Its total mass before normalization is $w_{v,b}$. These are exactly the two quantities needed in the grid estimate.

Expand \(x(h)y(k)\) using the bins in each component. For a term with a single bin \(b_v\) in each (take nonempty bins), its squared norm for the expectation on the right of \eqref{tra:collision-series:eq7}, including the indicator, is by \eqref{tra:collision-series:eq3} at most
\[
\begin{aligned}
&\left(\prod_v w_{v,b_v}\right)
\exp\left(-n+C n^{0.55}+\frac{C}{\log n}\sum_v \log r_{v,b_v}\right)\\
&\qquad\le
\exp(-n+C n^{0.55})\prod_v \left(C_1\, e^{-(A-C)b_v/\log n}\right).
\end{aligned}
\]
where \(C,C_1\) are absolute, and \(A\) can now be fixed greater than \(C+2\). Indeed \eqref{tra:collision-series:eq3} applies to the independent laws given by the squared kernels in the bins, normalized to probabilities. By triangle inequality on taking square roots for the sum over bins, geometric summation, and squaring again, we obtain
\[
\mathbb E_{H,K}\!\left[{\boldsymbol 1}_{hk\in KH}|x(h)y(k)|^2\right]
\ \le\
\exp(-n+C n^{0.55})\left(C_2\log n\right)^{2(l+m)}
\]
with another absolute \(C_2\). Finally,
\[
\log\left(\frac{|G|}{|H||K|}\right)=
\log\left(\frac{n!}{(m!)^l(l!)^m}\right)=n+O(\sqrt n\log n).
\]
Substitution in \eqref{tra:collision-series:eq7} proves \eqref{tra:collision-series:eq6}.
\end{proof}
\subsection{The positive tuple collision series}
\label{tra:collision-series:representation-normalization}
The preceding argument gave a weak bound for the complete regular trace. We next prove a separate low-level bound, keeping a weighted sum of positive single-copy traces throughout.

All sufficiently-large requirements here are absolute. This step will use standard symmetric group representation theory (parametrization by partitions, the branching rule, sign twist, hook length formula and regular representation multiplicities). For clarity, let \(\operatorname{Tr}_\lambda\) indicate a trace on a single copy of the irreducible indexed by \(\lambda\vdash n\), and \(f^\lambda\) its dimension, the number of standard tableaux. Our operators specified by group convolution can be considered in unitary representations throughout. Thus the trace in the regular representation has contributions \(f^\lambda \operatorname{Tr}_\lambda\).
\begin{lemma}\label{tra:collision-series:sparse-trace}
\label{tra:collision-series:sparse-trace-statement}
For small levels \(k\), by which we refer to \(k=n-\lambda_1\), we use the following estimate, independent of \eqref{tra:collision-series:eq6}, for all large enough \(n\):
\begin{equation}
B_k:=\sum_{\lambda:\ n-\lambda_1=k} f^{\bar\lambda}\operatorname{Tr}_\lambda \mathcal Q_n
\ \le\ n^{-0.01 k},
\qquad 1\le k\le n^{0.62}. \label{tra:collision-series:eq9}
\end{equation}
Here \(\bar\lambda\) is the partition obtained by omitting the first row. We include \(B_0\) in the notation in deriving the estimate, with empty tableau dimension 1.
\end{lemma}
\begin{proof}\label{tra:collision-series:sparse-trace-proof}
\label{tra:collision-series:collision-generating-function}
To prove \eqref{tra:collision-series:eq9}, let \(R_j^*\) in this paragraph and the derivation denote a number, namely the trace of \(\mathcal Q_n\) on the permutation representation on ordered \(j\)-tuples of distinct positions (including \(j=0\)). By branching, for \(j\le n/2\) we have
\[
R_j^*=\sum_{k=0}^j \binom jk B_k.
\]
Indeed the representation on tuples is induced from the trivial representation on a subgroup permuting \(n-j\) points. By Frobenius reciprocity and branching the multiplicity of \(\lambda\) is the number of standard tableaux for the skew diagram left by omitting a first row segment of length \(n-j\) (zero if the segment does not fit). Here if \(k\le j\) the skew part of the top row does not interact with \(\bar\lambda\) as \(j\le n/2\), giving the formula. Consequently
\begin{equation}
B_k=k!\,[z^k]e^{-z}\sum_{j=0}^n R_j^* z^j/j!, \quad k\le n/2. \label{tra:collision-series:eq10}
\end{equation}

There is a path expression for these tuple traces. Perform one pass on all the labels, and in its collision graph place an edge between two labels when they face each other at a switch. Labels may for this purpose be identified with their initial positions. This is a simple graph: once two paths face each other they differ in the bit produced there, which is not changed again, so they cannot face again. For \(s\ge 0\) let \(h_s\) for this graph be the number of edge subsets with exactly \(s\) vertices in their support (thus \(h_0=1\)). We claim the polynomial identity
\begin{equation}
\sum_{j=0}^n R_j^* z^j/j!
 =\mathbb E\sum_{s=0}^n h_s (z/n)^s(1+z/n)^{n-s}. \label{tra:collision-series:eq11}
\end{equation}
In computing the tuple trace, we sum squared probabilities for the \(j\)-tuple to move between prescribed positions over a pass, since its operator from the pass is multiplied by its adjoint. For a tuple start, sample a finish by using the switches. Given this start and finish, the trajectory of every marked label is determined, because each bit is used exactly once. The probability of the sampled tuple finish from the given start is
\[
n^{-j} 2^{\,\#\{\text{edges between marked labels}\}}.
\]
This follows since each marked path determines a bit at each layer, each a \(1/2\) factor, except that two facing paths use the same coin. Compatibility at collisions holds since we took a sampled finish. Taking the mean of this probability exactly gives the sum of squared probabilities for that start. Now choose an ordered uniform tuple start, so the marked set \(M\) is a uniform \(j\)-set independent of the collision graph in the full shuffle, and there are \((n)_j\) tuple starts. Expanding \(2^{\,\#\{\text{edges in }M\}}\) by counting edge subsets, the expansion multiplies \(h_s\), in expectation over \(M\), by \((j)_s/(n)_s\) (interpreted as zero for \(s>j\)). Including factors \((n)_j n^{-j}\) and summing proves \eqref{tra:collision-series:eq11}.

For any realized graph the maximum degree is at most \(d\), and the number of induced edges within any set of \(u\) vertices is at most
\begin{equation}
\tfrac12 u\log_2 u. \label{tra:collision-series:eq12}
\end{equation}
For \eqref{tra:collision-series:eq12}, consider each layer, grouping these vertices by their common initial suffix of bits after the layer coordinate. The suffix and the current layer coordinate are not yet updated before that layer. Within each suffix group every collision then is between the two subgroups for the values of that coordinate, with number at most \(\min(a,b)\) if their sizes are \(a,b\). Use
\[
2\min(a,b)\le (a+b)\log_2(a+b)-a\log_2 a-b\log_2 b,
\]
by concavity of binary entropy, taking zero terms as zero. Summing these bounds telescopes through the groupings by suffix and proves \eqref{tra:collision-series:eq12}.

The exact generating identity has reduced the weighted representation trace to edge subsets in one collision graph. The induced-edge bound controls the connected edge subsets on each support. We use it next to estimate the degree-$k$ coefficient uniformly over the required range.

For coefficient \(k\) we may truncate \eqref{tra:collision-series:eq11} by omitting \(s>k\). Bound the resulting expression times \(e^{-z}\) on \(|z|=r=k n^{0.015}\). Here \(1\le k\le n^{0.62}\) and \(r<n/2\) for large \(n\). Its modulus is at most
\begin{equation}
\exp(O(r^2/n))\ \mathbb E\sum_{s=0}^k h_s \left(\frac{r}{n-r}\right)^s, \label{tra:collision-series:eq13}
\end{equation}
bounding \(e^{-z}(1+z/n)^n\) and \(|n+z|^{-1}\). To control the sum, the number of connected edge subsets of support size \(u\ge 2\), in any of the graphs above, is at most
\[
n\, d^{2u}\, u^{u/2}.
\]
Indeed the connected vertex sets can be found by walks of length \(2(u-1)\) through a spanning tree, and given the set use \eqref{tra:collision-series:eq12}. A general edge subset is a collection of connected edge subsets. Dropping disjointness and truncation on totals, but still using \(u\le k\) for individual components, shows that the sum inside the expectation in \eqref{tra:collision-series:eq13} is at most
\[
\exp\left(\sum_{u=2}^k n\left(d^2\sqrt u\,\frac{r}{n-r}\right)^u\right).
\]
The exponent here is \(o(k)\) uniformly. For \(2\le u<40\) the total is \(O(n(d^2 r/n)^2)=o(k)\), since \(d=O(\log n)\) and \(k\le n^{0.62}\); for \(u\ge 40\) the number inside the parentheses raised to \(u\), before taking that power, is at most \(n^{-0.05}\) for all sufficiently large \(n\), so the total for those \(u\) is \(o(1)\). Also \(r^2/n=o(k)\). By Cauchy's coefficient bound in \eqref{tra:collision-series:eq10}, we get \(B_k\le k! r^{-k}\exp(o(k))\), proving \eqref{tra:collision-series:eq9}.
\end{proof}
\subsection{Closing the regular moment}
\begin{theorem}\label{tra:collision-series:moment}
\label{tra:collision-series:main-statement}
There are uniformly bounded real exponents \(q_n\ge 300\), for the powers of two \(n\ge 2\), such that
\begin{equation}
\operatorname{Tr} \mathcal Q_n^{q_n}\ \le\ 1+\tfrac18. \label{tra:collision-series:eq1}
\end{equation}
An absolute integer number of sweeps therefore has chi-square divergence at most $1/8$.
\end{theorem}
\begin{proof}
\label{tra:collision-series:bounded-exponents-and-completion}
The small-level input to the common moment closure is the positive
mass $B_k$, rather than a separate bound on each block. We first
convert that mass into a bound for the low-level regular trace.

First, the total contribution in the regular trace of \(\mathcal Q_n^q\), for any \(q\ge 300\), from levels \(1\le k\le n^{0.62}\), tends to zero. Indeed each such representation occurs in the corresponding \(k\)-tuple representation, so \(f^\lambda\le n^k\). The sum of the single-copy positive traces of \(\mathcal Q_n\) at level \(k\) is at most \(B_k\), hence the sum of their traces for the power \(q\) is at most \(B_k^q\). Thus the contributions including multiplicities are at most \(\sum_{1\le k\le n^{0.62}} n^k B_k^q=o(1)\).

If the transposed partition has level $k\le n^{0.62}$,
the original diagram has height $n-k>n/2$ for large $n$, and its
block vanishes by Lemma~\ref{lem:annihilation}.

All remaining diagrams, other than the trivial representation, have dimension at least \(\exp(n^{0.60})\) for large \(n\). Here are details of this coarse dimension bound. Let \(b=\lfloor n^{0.62}\rfloor\); the numbers outside the first row and outside the first column both exceed \(b\). If neither the first row nor first column has length \(5b\) or greater, all hook lengths are at most \(10b\), so the hook length formula gives a sufficient lower bound. Otherwise transpose if needed, which does not change dimension, to suppose the first row has length at least \(5b\). Use a subdiagram with first row of length \(4b\) and \(b\) more cells below it, chosen forming a partition there. This is possible by selecting a subdiagram of that size below the top row; in particular those cells have width at most \(b\). Its tableau dimension is a lower bound for the full dimension by extension of tableaux. By filling the first \(b\) cells in the top row first, and then interleaving a standard filling order in the lower \(b\) cells with the remaining top row cells, we get at least \(\binom{4b}b\) standard tableaux. This suffices as well.

Apply Lemma~\ref{lem:balanced-moment-closure} with
$\nu=2$, $\varepsilon=1/8$ and $P=300$, since
$\operatorname{Tr}\mathcal Q_n^q=\operatorname{Tr}|T_n|^{2q}$.
The low class consists of the levels $1\le k\le n^{0.62}$ and
the annihilated blocks. Its contribution is at most $1/16$ for all
large $n$, uniformly for $q\ge300$. The dimension calculation above
and Proposition~\ref{tra:collision-series:weak-moment} supply
\[
 H_n=n^{0.60},\qquad E_n=n^{0.57},\qquad
 \alpha_d=1+A/\log n.
\]
With $\delta_d=1/\log n$,
\[
 (1+\delta_d)E_n-\delta_dH_n
 =(1+1/\log n)n^{0.57}-n^{0.60}/\log n\longrightarrow-\infty,
\]
and $\alpha_d(1+\delta_d)=1+O(1/d)$. All thresholds are
independent of the child exponent. The lemma therefore proves
\eqref{tra:collision-series:eq1} with bounded $q_n$.

Taking an integer $s\ge\sup_nq_n$, the Schatten-power estimate
\eqref{e:holder-power} bounds the nonconstant regular squared norm
of $T_n^s$ by $1/8$. For convolution this is the chi-square
divergence from uniform; Proposition~\ref{e:nonnormal-moment-conversion}
gives total variation less than $1/4$. Thus an absolute number of
sweeps suffices. Lemma~\ref{lem:support} gives the lower bound
$2d-O(1)$ in physical shuffle units.
\end{proof}
